\documentclass{rvcepaper}
\begin{document}

% ---------- HS261TA Principles of Management and Economics
\begin{iempaper}
\paperinfo{shownote=0,date={15\textsuperscript{th} April 2026},exam={CIE - I},maxmarks={10 + 50},sem={VI},prog={UG},duration={30 + 90 Min},
 title={Principles of Management and Economics},code={HS261TA}}
\iemhead
{\small Note:}\par\vspace{-1pt}
\hspace*{0.5cm}{\footnotesize 1. \hspace{0.1cm}Answer all the Questions.}\par\vspace{3pt}
\begin{cietable}{M}{BT}{CO}
\qpart{Part -- A}
\q{1.}{List the five important functions of management.}{2}{1}{1}
\q{2.}{Identify which statement represents vision and which represents mission:
\begin{romi}\item[i.] ``To be a globally recognized brand known for innovation.''\item[ii.] ``We design and deliver innovative products to meet customer needs.''\end{romi}}{2}{1}{1}
\q{3.}{Do you agree that coordination need not be treated as a separate function of management? Give a brief reason.}{2}{2}{1}
\q{4.}{Explain how economists respond to the criticism that people often make irrational decisions due to emotions, habits, or lack of information.}{2}{2}{3}
\q{5.}{Explain how the government is involved in the circular flow of income.}{2}{2}{3}
\qpart{Part -- B}
\q{1}{A manager at Swiggy handles team coordination, information sharing, and decision-making during peak hours. Apply Henry Mintzberg's managerial roles to illustrate how the manager performs these functions in this situation.}{10}{L3}{1}
\q{2}{In the launch of a new electric vehicle by Tata Motors, the project manager must coordinate teams, solve problems, and motivate employees. Apply the important managerial skills to show how the manager can successfully handle this project.}{10}{L3}{1}
\q{3}{Using a real-life example, explain the concept of circular flow of income to show the role of households, firms, government, and foreign sector.}{10}{L2}{3}
\q{4}{Explain the basic economic problem of scarcity and choice and describe how individuals and firms allocate limited resources.}{10}{L2}{3}
\q{5}{Describe the different types of economic systems and explain any two with examples.}{10}{L2}{3}
\end{cietable}
\end{iempaper}

% ---------- IM362IA Supply Chain Management
\begin{iempaper}
\paperinfo{shownote=0,date={13/04/2026},exam={CIE - I},maxmarks={10 + 50},sem={VI},prog={UG},duration={30 + 90 Min},
 title={Supply Chain Management},code={IM362IA}}
\iemhead
{\small Note:}\par\vspace{-1pt}
\hspace*{0.5cm}{\footnotesize 1. \hspace{0.1cm}Answer all the Questions.}\par\vspace{3pt}
\begin{cietable}{M}{BT}{CO}
\qpart{Part--A}
\q{1.}{State the objective of a supply chain.}{1}{1}{CO1}
\q{2.}{List the typical stages in a supply chain.}{2}{1}{CO1}
\q{3.}{\blank[2cm] cycle occurs between the manufacturer and the supplier.}{1}{1}{CO1}
\q{4.}{Write the full form of the following supply chain macro processes: CRM, ISCM, SRM}{2}{2}{CO1}
\q{5.}{What is a company's supply chain strategy?}{2}{1}{CO1}
\q{6.}{Illustrate the relationship between number of facilities and total logistics cost.}{2}{2}{CO2}
\qpart{Part--B}
\q{1}{Explain the impact of supply chain decision phases on the overall profitability and success of manufacturing supply chains.}{10}{3}{CO1}
\q{2}{Illustrate and discuss cycle view of an automobile component manufacturing organization.}{10}{3}{CO1}
\q{3}{List the three basic steps to achieving this strategic fit. Evaluate the first step.}{10}{2}{CO1}
\q{4}{Unipart, a manufacturer of auto parts, is considering two B2B marketplaces to purchase its MRO supplies. Both marketplaces offer a full line of supplies at very similar prices for products and shipping. Both provide similar service levels and lead times. However, their fee structures are quite different. The first marketplace, Parts4u.com, sells all of its products with a 5 percent commission tacked on top of the price of the product (not including shipping). AllMRO.com's pricing is based on a subscription fee of \$10 million that must be paid up front for a two-year period and a commission of 1 percent on each transaction's product price Unipart spends about \$150 million on MRO supplies each year, although this varies with their utilization. Next year will likely be a strong year, in which high utilization will keep MRO spending at \$150 million. However, there is a 25 percent chance that spending will drop by 10 percent. The second year, there is a 50 percent chance that the spending level will stay where it was in the first year and a 50 percent chance that it will drop by another 10 percent. Unipart uses a discount rate of 20 percent. Assume all costs are incurred at the beginning of each year (so Year 1 costs are incurred now and Year 2 costs are incurred in a year). Which B2B marketplace should Unipart buy its parts from?}{20}{3}{CO2}
\end{cietable}
\end{iempaper}

% ---------- IM363IA Ergonomics
\begin{iempaper}
\paperinfo{shownote=0,date={15\textsuperscript{th} April 2026},exam={CIE - I},maxmarks={10 + 50},sem={6\textsuperscript{th}},prog={UG},duration={30 + 90 Min},
 title={ERGONOMICS},code={IM363IA}}
\iemhead
{\small Note:}\par\vspace{-1pt}
\hspace*{0.5cm}{\footnotesize 1. \hspace{0.1cm}Answer all the Questions.}\par\vspace{3pt}
\begin{cietable}{M}{BT}{CO}
\qpart{Part -- A}
\q{1.}{Define the term ergonomics.}{01}{L$_1$}{CO1}
\q{2.}{Name any two benefits of ergonomics}{01}{L$_2$}{CO1}
\q{3.}{What is meant by human centred design?}{01}{L$_1$}{CO1}
\q{4.}{Mention one example of cognitive ergonomics while interfacing with websites?.}{01}{L$_2$}{CO1}
\q{5.}{List Give a practical example of a Functional Anthropometric Data.}{01}{L$_2$}{CO2}
\q{6.}{What are the two important landmark percentiles in the normal distribution are as applied to consideration of anthropometric variables in designs?}{01}{L$_2$}{CO2}
\q{7.}{State any one principle of workspace design for sitting tasks.}{01}{L$_2$}{CO2}
\q{8.}{What is the primary goal of anthropometry?}{01}{L$_2$}{CO2}
\q{9.}{Give an example of application of Newtonian anthropometrics in designs}{01}{L$_3$}{CO2}
\q{10.}{Expand the acronyms CTD.}{01}{L$_1$}{CO2}
\qpart{Part -- B}
\q{1.}{Develop a step-by-step ergonomic evaluation plan for a computer based VDT based workstation. Your answer should include: Identification of key task elements, Application of anthropometric measurements, Postural analysis methods, Recommendations to reduce musculoskeletal risks.}{10}{L$_4$}{CO1}
\q{2.}{Distinguish between Macro ergonomics and Micro ergonomics. Outline the scope of work that is covered in both these areas of ergonomics. Give an example to justify the application of these specific to work systems design and equipment/workspace designs. List any two factors to be considered for each of these two types of ergonomics in designs.}{10}{L$_2$}{CO1}
\q{3.}{Consider the example of a handle Design of a Handheld Power Tool: How should you calculate the Grip diameter based on the 5\textsuperscript{th} and 95\textsuperscript{th} percentile values of a hand breadth. For easy - Thumb reach: where should the controls be placed considering natural thumb arc. Which force can be reduced by choosing - Material: such as soft rubberized textures? Which of the two grips whether inline grip or pistol grip must be so as to minimize wrist deviation. If you can reduce the handle mass to $\leq$ 1 kg and achieve the balance centre near grip, what is the implications for operator fatigue during the power tool operations? (Provide Sketches to highlight the engineering anthropometric design and to visualize the handles design of a Handheld power tool)}{10}{L$_4$}{CO2}
\q{4.}{A company is designing a sitting assembly station for both male and female operators in a watch assembly line. Apply anthropometric principles to determine optimal work surface height, Analyse how human variability influences design decisions, Discuss trade-offs between productivity, comfort, and injury prevention in this work context. (Sketch your ideas with clarity of critical design variables)}{10}{L$_3$}{CO2}
\q{5a.}{What is ``dynamic sitting,'' and why is it preferred over maintaining a ``perfect'' upright posture?}{05}{L$_2$}{CO1}
\q{b.}{The term Maximum Holding Time (MHT) is defined as The longest duration a specific physical posture or force can be maintained before failure or unacceptable discomfort. Give a practical example to emphasize as to why this variable is important for equipment/workspace design based on your own exposure and experience.}{05}{L$_3$}{CO2}
\end{cietable}
\end{iempaper}

% ---------- IM364TA Human Resource Management & Analytics
\begin{iempaper}
\paperinfo{shownote=0,date={16\textsuperscript{th} April 2026},exam={CIE - I},maxmarks={10 + 50},sem={VI},prog={UG},duration={30 + 90 Min},
 title={Human Resource Management \& Analytics},code={IM364TA}}
\iemhead
{\small Note:}\par\vspace{-1pt}
\hspace*{0.5cm}{\footnotesize Answer all the Questions.}\par\vspace{3pt}
\begin{cietable}{M}{BT}{CO}
\qpart{Part -- A}
\q{1.}{HR manager advises a department head on hiring. What type of authority is this?}{01}{L1}{CO1}
\q{2.}{A company uses online interviews for hiring. Which trend is this in HRM?}{01}{L2}{CO1}
\q{3.}{Which HR group focuses on strategic decision-making?}{01}{L1}{CO1}
\q{4.}{What is talent management?}{01}{L1}{CO1}
\q{5.}{What are the key components of a job description?}{01}{L1}{CO1}
\q{6.}{What is HR forecasting?}{01}{L1}{CO2}
\q{7.}{List the three steps in Succession planning?}{02}{L2}{CO2}
\q{8.}{List any two internal sources of candidates.}{01}{L1}{CO2}
\q{9.}{Define Job analysis}{01}{L1}{CO1}
\qpart{Part -- B}
\q{1.}{Explain the major trends shaping Human Resource Management in detail.}{10}{L2}{CO1}
\q{2.}{Imagine you are working in the HR department of a hospital called \textit{CarePlus Medical Center}. The hospital wants to improve its hiring process for the role of Staff Nurse. Prepare:\newline
a) A Job Description for the role of Staff Nurse\newline
b) A Job Specification for the same role}{10}{L4}{CO1}
\q{3.}{A manufacturing company wants to analyze a job that involves repetitive manual tasks on the shop floor. Which job analysis methods would you recommend? Justify your answer by comparing alternative methods.}{10}{L3}{CO1}
\q{4.}{An e-commerce company is expanding and needs to hire for marketing, customer support, data analysis, and logistics roles. Identify different external sources for recruiting candidates and explain how each aligns with the company's needs.}{10}{L3}{CO2}
\q{5.}{Imagine you are the HR manager of a multinational company that is about to launch a recruitment campaign for a new product development team. The team will consist of software developers, project managers, and quality assurance specialists.\newline
a) Outline the steps of the recruitment process you would follow to attract and hire the best talent for this team.\newline
b) Explain how you would ensure that the process is efficient and cost-effective.\newline
Discuss any tools or technologies you might use to streamline this recruitment process.}{10}{L4}{CO2}
\end{cietable}
\end{iempaper}

% ---------- IM365TDA Facilities Planning and Design
\begin{iempaper}
\paperinfo{shownote=0,date={15\textsuperscript{th} April 2026},exam={CIE - I},maxmarks={10 + 50},sem={VI},prog={UG},duration={30 + 90 Min},
 title={FACILITIES PLANNING AND DESIGN},code={IM365TDA}}
\iemhead
{\small Note:}\par\vspace{-1pt}
\hspace*{0.5cm}{\footnotesize 1. \hspace{0.1cm}Answer all the Questions.}\par\vspace{3pt}
\ptitle{PART -- A}
\begin{longtable}{|C{1.05cm}|p{12.2cm}|C{0.95cm}|C{0.95cm}|C{0.95cm}|}
\hline \textbf{SI.\newline No.} & \multicolumn{1}{c|}{\textbf{Questions}} & \textbf{M} & \textbf{BT} & \textbf{CO}\\\hline
\q{1}{State any two factors influencing plant location.}{02}{L2}{CO1}
\q{2}{Identify the suitable layout for a hospital and justify your choice.}{02}{L2}{CO1}
\q{3}{Which layout is suitable for job-order production? Give reason.}{02}{L1}{CO2}
\q{4}{Analyze any two causes of inadequate facilities planning with examples.}{02}{L1}{CO1}
\q{5}{Propose a basic facilities planning strategy for a new manufacturing firm.}{02}{L2}{CO2}
\end{longtable}
\ptitle{PART -- B}
\begin{longtable}{|C{1.05cm}|p{12.2cm}|C{0.95cm}|C{0.95cm}|C{0.95cm}|}
\hline \textbf{SI.\newline No.} & \multicolumn{1}{c|}{\textbf{Questions}} & \textbf{M} & \textbf{BT} & \textbf{CO}\\\hline
\q{1.}{Explain the concept of facilities planning and discuss its significance in modern industries.}{10}{L1}{CO1}
\q{2.}{Differentiate between facilities planning process and strategic planning process in terms of meaning, scope, objective, time horizon and focus.}{10}{L2}{CO2}
\q{3.}{Explain the factors influencing plant location with suitable examples.}{10}{L3}{CO3}
\q{4.}{Recommend a suitable plant layout for a manufacturing firm and justify your choice with merits and demerits.}{10}{L4}{CO4}
\q{5.}{Potential locations A, B and C have the cost structure as shown below for producing a product expected to sell for Rs. 4250 per unit.
\qtabc{|c|c|c|}{\hline Site & Fixed Cost/year (Rs) & Variable Cost/unit (Rs)\\\hline A & 85,00,000 & 775\\\hline B & 65,00,000 & 2225\\\hline C & 70,00,000 & 1460\\\hline}
Find the most economical location for an expected volume of 6250 units per year.}{10}{L2}{CO2}
\end{longtable}
\end{iempaper}

% ---------- IM365TDB Service Operations Management
\begin{iempaper}
\paperinfo{shownote=0,date={16\textsuperscript{th} April 2026},exam={CIE - 1},maxmarks={10 + 50},sem={VI},prog={UG},duration={30 + 90 Min},
 title={Service Operations Management},code={IM365TDB}}
\iemhead
{\small Note:}\par\vspace{-1pt}
\hspace*{0.5cm}{\footnotesize 1. \hspace{0.1cm}Answer all the Questions.}\par\vspace{3pt}
\begin{cietable}{M}{BT}{CO}
\qpart{Part -- A}
\q{1.}{What is the primary goal of Service Operations Management?}{2}{1}{1}
\q{2.}{Name two major challenges faced by service operations managers in today's dynamic environment.}{2}{2}{2}
\q{3.}{Identify any 4 tools that have been used on you to collect information to enhance customer satisfaction.}{2}{3}{1}
\q{4.}{What type of service process is typically used in a fast-food restaurant, and why?}{2}{3}{2}
\q{5.}{How does service operations management differ from manufacturing operations management?}{2}{2}{2}
\qpart{Part -- B}
\q{1}{What are the different types of service processes? Cite the example of any service organization and highlight the different processes involved in that organization.}{10}{2}{2}
\q{2}{Apply the concept of the ``Poka-Yoke'' (mistake-proofing) methodology to a service setting of your choice. Explain how you would design both warning and control types of poka-yokes to manage the operational risks associated with customer variability and employee variability.}{10}{3}{2}
\q{3}{Explain ``Service concept'' using the tabular column as applied to a ``New age engineering education institute''.}{10}{5}{3}
\q{4}{You are tasked with launching a new, subscription-based ``tele-health'' service that connects patients with doctors via video calls 24/7. This service is highly intangible and faces significant variability in demand.\newline
Create a set of success metrics that specifically measure the ``moments of truth'' in the virtual interaction, addressing the challenge of measuring quality for an intangible product.\newline
Hint: Go up to 3 Moments of Truth starting from Zeroth. Use tabular column and avoid writing in paragraphs.}{10}{4}{2}
\q{5}{You are a Service Operations Manager of an online retail store. What challenges would you may face as the manager and what measures would you take to overcome them?}{10}{3}{3}
\end{cietable}
\end{iempaper}

% ---------- IM365TDD Design of Experiments
\begin{iempaper}
\paperinfo{shownote=0,date={15\textsuperscript{th} April 2026},exam={CIE - I},maxmarks={10 + 50},sem={VI},prog={UG},duration={30 + 90 Min},
 title={Design of Experiments},code={IM365TDD}}
\iemhead
{\small Note:}\par\vspace{-1pt}
\hspace*{0.5cm}{\footnotesize 1. \hspace{0.1cm}Answer all the Questions.}\par
\hspace*{0.5cm}{\footnotesize 2. \hspace{0.1cm}Use of statistical tables permitted}\par\vspace{3pt}
\begin{cietable}{M}{BT}{CO}
\qpart{Part -- A}
\q{1.}{List two applications of design of experiments.}{1}{1}{CO1}
\q{2.}{What are the two advantages of replication?}{2}{1}{CO1}
\q{3.}{\blank[2.4cm] is an example of larger-the-better type quality characteristic.}{1}{1}{CO2}
\q{4.}{According to Genichi Taguchi ``Quality is the \blank[1.8cm] a product causes to the society''}{2}{2}{CO2}
\q{5.}{Identify response, controllable and noise factors for gas welding process.}{2}{2}{CO1}
\q{6.}{What are the steps in the engineering design problem strategy?}{2}{2}{CO2}
\qpart{Part -- B}
\q{1}{Enumerate the guidelines for the statistical approach in designing and analyzing an experiment.}{10}{2}{CO1}
\q{2}{Trace the history of statistical design. Elucidate the important stages.}{10}{3}{CO1}
\q{3}{Write the expressions for quality loss function of various quality characteristics. Illustrate with sketches. List the situations in which these quality characteristics are applicable.}{10}{2}{CO2}
\q{4}{Identify the quality control activities during various product realization steps. Comment on the ability of these activities to reduce effect of noise factors.}{05}{3}{CO2}
\q{5}{An industrial engineer employed by a beverage bottler is interested in the effects of two different types of 32-ounce bottles (A) on the time to deliver 12-bottle cases of the product. The two bottle types are glass and plastic. Two workers (B) are used to perform a task consisting of moving 40 cases of the product 50 feet on a standard type of hand truck and stacking the cases in a display. Four replicates of a $2^2$ factorial design are performed, and the times observed are listed in the following table:
\qtab{|c|c|c|c|c|}{\hline Treatment & \multicolumn{4}{c|}{Replicates}\\\cline{2-5} Combination & I & II & III & IV\\\hline 1 & 5.12 & 4.89 & 4.98 & 5.00\\\hline a & 6.65 & 6.24 & 5.49 & 5.55\\\hline b & 4.95 & 4.43 & 4.27 & 4.25\\\hline ab & 5.28 & 4.91 & 4.75 & 4.71\\\hline}
\begin{romi}\item Estimate the factor effects.\item Prepare an analysis of variance table and determine which factors are important in explaining time to deliver.\end{romi}}{15}{2}{CO4}
\end{cietable}
\end{iempaper}

% ---------- IM266TEQ Elements of Financial Management
\begin{iempaperB}
\paperinfo{hdr=2,ayear={2025-2026 (EVEN Sem)},date={16\textsuperscript{th} April 2026},code={IM266TEQ},sem={VI Semester},exam={CIE - I},maxmarks={10 + 50},duration={120 Mins},
 title={ELEMENTS OF FINANCIAL MANAGEMENT (Institutional Elective)},shownote=0}
\iemheadB
{\small Note:}\par\vspace{-1pt}
\hspace*{0.5cm}{\footnotesize 1. \hspace{0.1cm}Answer all the Questions.}\par
\hspace*{0.5cm}{\footnotesize 2. \hspace{0.1cm}Use of Interest factor tables permitted.}\par\vspace{3pt}
\ptitle{PART -- A}
\begin{xtable}{Questions}
\q{1}{Distinguish between Financial Assets and Liabilities.}{2}{2}{2}
\q{2}{Mention the role of SEBI in investor protection.}{2}{2}{2}
\q{3}{What is meant by the term, ``Bottom Line'' in the context of financial management?}{2}{2}{1}
\q{4}{What is meant by the term Depreciation? Mention how it is treated in preparation of Accounts.}{2}{2}{1}
\q{5}{Distinguish between ordinary annuity and annuity due.}{2}{2}{2}
\end{xtable}
\ptitle{PART -- B}
\begin{xtable}{Questions}
\q{1.}{Analyze the conflict between profit maximization and shareholder wealth maximization in modern corporations.}{10}{4}{1}
\q{2.}{Assess the effectiveness of India's regulatory framework (RBI, SEBI, Companies Act) in ensuring financial stability.}{10}{4}{1}
\q{3.}{Evaluate the role of financial intermediaries in promoting financial inclusion in India.}{10}{5}{1}
\q{4.a}{Compare the usefulness of balance sheet vs profit and loss statement for investors.}{05}{2}{2}
\q{b}{Compute Future Value of an amount \rupee 1,000 invested at 8\% interest rate for a period of 5 years.}{05}{3}{2}
\q{5.a}{Compute Future Value of an amount \rupee 2,000 invested annually for a period of 6 years at 12\% rate of interest.}{05}{3}{2}
\q{b}{A tech start-up is promised a \$1M pay out in 10 years. However, inflation is expected to be 3\% for the first 4 years and 5\% for the remaining 6 years. Calculate the PV of this pay out in today's dollars.}{05}{3}{2}
\end{xtable}
\stars{***********}
\end{iempaperB}

% ---------- HS266TEY Universal Human Values - III
\begin{xpaper}
\paperinfo{deptline={},date={16th Apr 2026},exam={CIE -- I (Test)},maxmarks={50+10},sem={VI},prog={},duration={90+20 Min},
 title={Universal Human Values -- III},code={HS266TEY}}
\noindent\begingroup\renewcommand{\arraystretch}{1.5}\bfseries\fontsize{11}{13}\selectfont
\begin{tabular}{|p{3.2cm}|p{4.8cm}|C{3.6cm}|C{3.6cm}|}\hline
 Date & 16th Apr 2026 & Maximum Marks & 50+10\\\hline
 Course Code & HS266TEY & Duration & 90+20 Min\\\hline
 Sem & VI & \multicolumn{2}{l|}{CIE -- I (Test)}\\\hline
 \multicolumn{4}{|c|}{Universal Human Values -- III}\\\hline
\end{tabular}\endgroup\par\vspace{8pt}
\begin{longtable}{|C{1.05cm}|p{12.2cm}|C{0.95cm}|C{0.95cm}|C{0.95cm}|}
\hline \textbf{Sl. No.} & \multicolumn{1}{c|}{\textbf{Quiz}} & \textbf{M} & \textbf{BT} & \textbf{CO}\\\hline
\q{1}{\blank[1.6cm] is central to human existence, \blank[1.6cm] is central to \blank[1.6cm] existence.}{02}{1}{2}
\q{2}{Define co-existence.}{02}{2}{1}
\q{3}{The basic human aspiration is fulfilled through \blank[1.6cm] \& \blank[1.6cm]\ .}{02}{1}{1}
\q{4}{\blank[1.6cm] from animal consciousness to human consciousness is ensured through \blank[1.6cm].}{02}{2}{2}
\q{5}{Living with human consciousness provides base for \blank[1.6cm] \& \blank[1.6cm] ensuring justice and order leading to \blank[1.6cm].}{02}{1}{1}
\end{longtable}
\par\vspace{6pt}
\begin{longtable}{|C{1.05cm}|p{12.2cm}|C{0.95cm}|C{0.95cm}|C{0.95cm}|}
\hline \textbf{Sl. No.} & \multicolumn{1}{c|}{\textbf{Test}} & \textbf{M} & \textbf{BT} & \textbf{CO}\\\hline
\q{1}{Justify the statement ``Human Consciousness provides base for Universal Human Order''.}{10}{1}{1}
\q{2}{Explain the concept of basic human aspiration and its fulfillment.}{10}{2}{2}
\q{3}{Describe Continuous Happiness as the Basic Human Aspiration.}{10}{3}{2}
\q{4}{Elaborate on personal and societal transformation.}{10}{3}{3}
\q{5}{Discuss the process of self-exploration in detail.}{10}{2}{2}
\end{longtable}
\par\vspace{4pt}\noindent{\bfseries Course outcomes:}\par
\noindent\begin{tabular}{|C{1.3cm}|p{15cm}|}\hline CO 1 & Understand the basic human aspiration with program of its fulfillment and meaning of resolution in the complete expanse of human living\\\hline CO 2 & Understand human being in depth and see how self is central to human being\\\hline CO 3 & Understand existence in depth and see how coexistence is central to existence\\\hline CO 4 & Understand human conduct and the holistic way of living leading to human tradition\\\hline\end{tabular}
\btfoot{|l|l|l|*{10}{C{0.9cm}|}}{\hline \multirow{2}{*}{Marks Distribution} & Particulars & & CO1 & CO2 & CO3 & CO4 & L1 & L2 & L3 & L4 & L5 & L6\\\cline{2-13} & Test & Max Marks & 16 & 34 & 10 & -- & 16 & 24 & 20 & -- & -- & --\\\hline}
\stars{**********}
\end{xpaper}

% ---------- HS266TEW Applied Psychology for Engineers
\begin{iempaper}
\paperinfo{shownote=0,date={16.04.2026},exam={CIE - I},maxmarks={10 + 50},sem={VI},prog={UG- Institutional Elective},duration={2 hrs},
 title={Applied Psychology for Engineers},code={HS266TEW}}
\begin{center}{\fontsize{12}{14}\selectfont DEPARTMENT OF}\\[3pt]{\bfseries\fontsize{15.5}{18}\selectfont HUMANITIES AND SOCIAL SCIENCES}\end{center}\vspace{-3pt}
\noindent\begingroup\renewcommand{\arraystretch}{2.0}\bfseries\fontsize{11.5}{13}\selectfont
\begin{tabular}{|p{6.2cm}|C{5.0cm}|p{5.5cm}|}\hline
 Date \hspace{0.5cm}: \hspace{0.2cm}16.04.2026 & CIE - I & Max. Marks \hspace{0.3cm}: \hspace{0.2cm}10 + 50\\\hline
 Semester \hspace{0.1cm}: \hspace{0.2cm}VI & UG- Institutional Elective & Duration \hspace{0.55cm}: \hspace{0.2cm}2 hrs\\\hline
 \multicolumn{2}{|p{11.2cm}|}{Course Title: Applied Psychology for Engineers} & Course Code \hspace{0.1cm}: \hspace{0.2cm}HS266TEW\\\hline
\end{tabular}\endgroup\par\vspace{4pt}
\begin{cietable}{M}{BT}{CO}
\qpart{Part A}
\q{1.1}{Who is known as the founder of behaviourism in Applied Psychology?}{1}{L1}{CO1}
\q{1.2}{Identify the subfield of applied psychology in the following scenario: A company hires a specialist to improve employee motivation and productivity in the workplace.}{1}{L1}{CO1}
\q{1.3}{Which branch of psychology examines how peer pressure influences individual behaviour in groups?}{1}{L1}{CO2}
\q{1.4}{In which research method are variables systematically manipulated to observe their effect on behaviour?}{1}{L2}{CO2}
\q{1.5}{State the main objectives of applied psychology in real-world settings.}{2}{L2}{CO2}
\q{1.6}{Which field of psychology focuses on assessing and treating emotional and behavioural disorders?}{1}{L1}{CO2}
\q{1.7}{Rahul is facing anxiety and stress related to his academic performance. Which type of psychologist should he consult?}{1}{L3}{CO4}
\q{1.8}{Who proposed the concept of primary mental abilities?}{1}{L1}{CO3}
\q{1.9}{Which ability helps in predicting future performance?}{1}{L1}{CO2}
\multicolumn{5}{|c|}{50 Marks}\\ \hline
\qpart{Part B}
\q{2}{``Sadie believes that she is the Queen of England and that the CIA is out to get her''. Explain in detail the type of psychologist would be most likely to treat her psychological disorder?}{10}{L1}{CO1}
\q{3}{Explain in detail the approach of psychology that emphasizes conscious experiences, including each person's unique potential for psychological growth and self-direction.}{10}{L2}{CO2}
\q{4}{Dr. Stephen is going to observe and record children's play behaviour at a nursery school without attempting to influence or control the behaviour. Which method of research is involved? Explain the process and discuss its merits and demerits.}{10}{L3}{CO2}
\q{5}{Dev shows consistent performance across verbal, numerical, and spatial tasks. How would Spearman interpret this consistency? Justify the statement.}{10}{L1}{CO1}
\q{6}{A student shows high ability in memory and word fluency but average performance in reasoning tasks. Using Thurston's model, justify this pattern of abilities.}{10}{L2}{CO2}
\end{cietable}
\end{iempaper}

% ---------- AS266TEA Fundamentals of Aerospace Engineering (CIE-1)
\begin{iempaperB}
\paperinfo{hdr=2,usn=0,rule=0,ayear={2025-2026 (Even Semester)},dept={AEROSPACE ENGINEERING},date={13\textsuperscript{th} April 2024},code={AS266TEA},sem={VI Semester},exam={CIE-1},maxmarks={60},duration={120 Min},
 title={FUNDAMENTALS OF AEROSPACE ENGINEERING},shownote=0}
\iemheadB
{\small Note:\ \ 1. Answer All Questions from both Part-A \& Part-B}\par
\hspace*{0.9cm}{\small 2. Part-A to be answered in the 1\textsuperscript{st} two pages of the answer booklet}\par\vspace{4pt}
\ptitlem{PART -- A}{(10 Marks)}
\begin{longtable}{|C{1.0cm}|p{11.7cm}|C{1.2cm}|C{1.2cm}|C{1.2cm}|}
\hline & \multicolumn{1}{c|}{\textbf{QUESTIONS}} & \textbf{M} & \textbf{CO} & \textbf{BL}\\\hline
\q{1.}{List the ISA relevance in design and Performance of aircraft}{02}{CO1}{L2}
\q{2.}{The main source of thrust on large commercial airliners is/are \blank[2cm].}{01}{CO1}{L2}
\q{3.}{The pitching motion/moment in an aircraft is controlled by?}{01}{CO1}{L2}
\q{4.}{What are the three primary control surfaces in an aircraft.}{01}{CO2}{L2}
\q{5.}{Calculate the geopotential altitude(h) if $r = 6.356810^{6}$m and $h_G$ = 7 Km}{01}{CO2}{L2}
\q{6.}{The vertical straight line in the temperature distribution of standard atmosphere graph is called as \blank[2cm] region}{01}{CO1}{L1}
\q{7.}{What is the standard sea level value of pressure, temperature and density?}{01}{CO3}{L2}
\q{8.}{The rudder is located on the \blank[2cm] of aircraft}{01}{CO1}{L1}
\q{9.}{Why is the International Standard Atmosphere (ISA) considered a Hypothetical model?}{01}{CO2}{L1}
\end{longtable}
\newpage
\ptitle{PART-B}
\begin{longtable}{|C{1.0cm}|p{11.7cm}|C{1.2cm}|C{1.2cm}|C{1.2cm}|}
\hline \textbf{Sl.\newline No.} & \multicolumn{1}{c|}{\textbf{Questions}} & \textbf{M} & \textbf{CO} & \textbf{BT}\\\hline
\q{1}{At what value of geometric altitude $h_G$ is the difference $h$- $h_G$ equal to 2\% of $h$?}{10}{CO2}{L3}
\q{2}{Identify the major components and its function of a typical Helicopter with a neat sketch.}{10}{CO3}{L2}
\q{3}{Construct the expression to relate the P, T and $\rho$ along the gradient and isothermal regions of the atmosphere.}{10}{CO3}{L1}
\q{4}{Calculate the Standard Atmospheric values of P, T and $\rho$ at a geo-potential altitude of 14km.}{10}{CO1}{L2}
\q{5}{List the various altitudes with a neat sketch and establish a relation expressing the dependence of geometric and geo-potential altitudes.}{10}{CO4}{L1}
\end{longtable}
\btfoot{|l|l|l|*{10}{C{0.9cm}|}}{\hline
\multirow{2}{*}{\shortstack{Marks\\Distribution}} & Particulars & & CO1 & CO2 & CO3 & CO4 & L1 & L2 & L3 & L4 & L5 & L6\\\cline{2-13}
 & Test & \shortstack{Max\\Marks} & 16 & 13 & 21 & 10 & 21 & 29 & 10 & XX & XX & XX\\\hline}
\stars{***********}
\end{iempaperB}

\end{document}
